% Part: first-order-logic
% Chapter: models-theories
% Section: expressing-relations

\documentclass[../../../include/open-logic-section]{subfiles}

\begin{document}

\olfileid{fol}{mat}{exr}

\olsection{\article{structure}
  \printtoken{S}{structure}లో సంబంధాలను వ్యక్తం చేయడం}

\begin{explain}
\tetoken{సూత్రాలు}{!!{formula}s} యొక్క ఒక ప్రధాన ఉపయోగం, \tetoken{ఒక నిర్మాణం}{!!a{structure}}~$\Struct M$లోని
ధర్మాలు, సంబంధాలను~$\Struct M$ భాష~$\Lang L$ యొక్క ఆదిమ భావనల ద్వారా
వ్యక్తం చేయడం. దీని అర్థం ఇలా: $\Struct M$ యొక్క \tetoken{వ్యక్తి క్షేత్రం}{!!{domain}} కొన్ని
వస్తువుల సమితి. \tetoken{స్థిరసంకేతాలు}{!!{constant}s}, \tetoken{ప్రమేయ సంకేతాలు}{!!{function}s}, \tetoken{విధేయ సంకేతాలకు}{!!{predicate}s}
$\Struct M$లో వరుసగా~$\Domain M$లోని కొన్ని వస్తువులు, $\Domain M$పై
ప్రమేయాలు, $\Domain M$పై సంబంధాలు అర్థనిర్దేశాలుగా ఉంటాయి. ఉదాహరణకు,
$\Obj A^2_0$, $\Lang L$లో ఉంటే, $\Struct M$ దానికి
$R = \Assign{{\Obj A^2_0}}{M}$ అనే సంబంధాన్ని కేటాయిస్తుంది. అప్పుడు
సూత్రం~$\Atom{\Obj A^2_0}{\Obj v_1, \Obj v_2}$ కింది అర్థంలో సరిగ్గా ఆ
సంబంధాన్ని \emph{వ్యక్తం చేస్తుంది}: చర నిర్దేశం~$s$, $\Obj v_1$ను
$a \in \Domain{M}$కు, $\Obj v_2$ను $b \in \Domain M$కు చిత్రిస్తే,
\[
Rab \text{\quad iff\quad} \Sat{M}{\Atom{\Obj A^2_0}{\Obj v_1, \Obj v_2}}[s].
\]
ఇక్కడ చర నిర్దేశాలను చేర్చడం అవసరమని గమనించండి: ``$Rab$ అయితేనే
$\Sat{M}{\Atom{\Obj A^2_0}{a, b}}$'' అని మాత్రమే చెప్పలేం. ఎందుకంటే
$a$, $b$ మన భాషలోని సంకేతాలు కావు; అవి~$\Domain{M}$లోని
\tetoken{మూలకాలు}{!!{element}s}.

మనకు పరమాణు \tetoken{సూత్రాలు}{!!{formula}s} మాత్రమే కాకుండా, తార్కిక సంయోజకాలు,
పరిమాణీకరణికాల ద్వారా వాటిని కలపగల అవకాశం కూడా ఉంది కాబట్టి, మరింత
సంక్లిష్టమైన \tetoken{సూత్రాలు}{!!{formula}s} నేరుగా~$\Struct M$లో నిర్మితమై లేని ఇతర
సంబంధాలను నిర్వచించగలవు. అది ఎలా చేయాలో, ప్రత్యేకంగా
\tetoken{ఒక నిర్మాణంలో}{!!a{structure}} ఏ
సంబంధాలను నిర్వచించగలమో తెలుసుకోవడం మన ఆసక్తి.
\end{explain}

\begin{defn}
$!A(\Obj v_1,\dots, \Obj v_n)$ అనేది~$\Lang L$లోని \tetoken{ఒక సూత్రం}{!!a{formula}} అని,
దానిలో $\Obj v_1$,\dots, $\Obj v_n$ మాత్రమే స్వేచ్ఛగా సంభవిస్తాయని,
$\Struct M$ అనేది~$\Lang L$ కోసం \tetoken{ఒక నిర్మాణం}{!!a{structure}} అని అనుకుందాం. ఏ చర
నిర్దేశం~$s$కైనా $s(\Obj v_i) = a_i$ ($i = 1, \dots, n$) అయినప్పుడు
\[
Ra_1\dots a_n \text{\quad iff\quad} \Sat{M}{\Atom{!A}{\Obj
    v_1,\dots, \Obj v_n}}[s]
\]
అయితే, అప్పుడు మాత్రమే $!A(\Obj v_1,\dots, \Obj v_n)$,
$R \subseteq \Domain M^n$ అనే \emph{సంబంధాన్ని వ్యక్తం చేస్తుంది}.
\end{defn}

\begin{ex}
అంకగణిత ప్రామాణిక నమూనా~$\Struct N$లో, \tetoken{సూత్రం}{!!{formula}}
$\Obj v_1 < \Obj v_2 \lor \eq[\Obj v_1][\Obj v_2]$, $\Nat$పై
$\le$ సంబంధాన్ని వ్యక్తం చేస్తుంది. \tetoken{సూత్రం}{!!{formula}}
$\eq[\Obj v_2][\Obj v_1']$ ఉత్తరవర్తి సంబంధాన్ని, అంటే $m$ అనేది
$n$కు ఉత్తరవర్తి అయినప్పుడు $Rnm$ నిలిచే $R \subseteq \Nat^2$
సంబంధాన్ని వ్యక్తం చేస్తుంది. సూత్రం~$\eq[\Obj v_1][\Obj v_2']$
పూర్వవర్తి సంబంధాన్ని వ్యక్తం చేస్తుంది. \tetoken{సూత్రాలు}{!!{formula}s}
$\lexists[\Obj v_3][(\eq/[\Obj v_3][\Obj 0] \land
  \eq[\Obj v_2][(\Obj v_1 + \Obj v_3)])]$ మరియు $\lexists[\Obj
  v_3][\eq[(\Obj v_1 + {\Obj v_3}')][\Obj v_2]]$ రెండూ~$\Obj <$
సంబంధాన్ని వ్యక్తం చేస్తాయి. అంటే అంకగణిత భాషలో విధేయ సంకేతం~$<$
నిజానికి అనవసరమే; దానిని నిర్వచించవచ్చు.
\sourcecorrection{OLTEFOLMAT-001}{రెండవ సూత్రంలో మూలం చివరి చరాన్ని $v_2$గా మాత్రమే ముద్రించింది; అదే సూత్రంలోని ఇతర వస్తుభాషా చరాలకు, సమాంతర ఉదాహరణకు అనుగుణంగా $\Obj v_2$గా సరిచేశాం.}
\end{ex}

\begin{explain}
ఈ భావన నిర్దిష్ట \tetoken{నిర్మాణాలలో}{!!{structure}s} మాత్రమే ఆసక్తికరమైనది కాదు; ఒక
ఉద్దేశించిన నమూనాను లేదా నమూనాలను వర్ణించడానికి ఒక భాషను ఉపయోగించే
ప్రతి సందర్భంలోనూ, అంటే సిద్ధాంతాలను పరిశీలించినప్పుడల్లా, ఇది
ఆసక్తికరమే. ఈ సిద్ధాంతాల్లో తరచుగా కొన్ని \tetoken{విధేయ సంకేతాలు}{!!{predicate}s} మాత్రమే
మౌలిక సంకేతాలుగా ఉంటాయి; కానీ అవి వర్ణించే క్షేత్రంలో అనేక ఇతర
సంబంధాలు ముఖ్యమైన పాత్ర పోషిస్తాయి. భాషలోని మౌలిక \tetoken{విధేయ సంకేతాలకు}{!!{predicate}s}
అర్థనిర్దేశాలైన సంబంధాల ద్వారా ఈ ఇతర సంబంధాలను క్రమబద్ధంగా వ్యక్తం
చేయగలిగితే, వాటిని ఆ భాషలో \emph{నిర్వచించగలం} అంటాం.
\end{explain}

\begin{prob}
కింది సంబంధాలను నిర్వచించే \tetoken{సూత్రాలను}{!!{formula}s}~$\Lang L_A$లో కనుగొనండి:
\begin{enumerate}
\item $n$ అనేది $i$, $j$ల మధ్య ఉంది;
\item $n$, $m$ను శేషం లేకుండా భాగిస్తుంది (అంటే, $m$ అనేది $n$కు
  ఒక గుణితం);
\item $n$ ఒక ప్రధాన సంఖ్య (అంటే, $1$, $n$ తప్ప మరే సంఖ్యా~$n$ను శేషం
  లేకుండా భాగించదు).
\end{enumerate}
\end{prob}

\begin{prob}
సూత్రం~$!A(\Obj v_1, \Obj v_2)$, \tetoken{ఒక నిర్మాణం}{!!a{structure}}~$\Struct M$లో
$R \subseteq \Domain M^2$ సంబంధాన్ని వ్యక్తం చేస్తుందనుకుందాం. కింది
సంబంధాలను వ్యక్తం చేసే సూత్రాలను కనుగొనండి:
\begin{enumerate}
\item $R$ యొక్క విలోమం~$R^{-1}$;
\item సాపేక్ష లబ్ధం~$R \mid R$.
\end{enumerate}
$R$ యొక్క సంక్రామక సంవృతి~$R^+$ను వ్యక్తం చేసే మార్గాన్ని
కనుగొనగలరా?
\end{prob}

\begin{prob}
$\Lang{L}$లో $<$ అనే ద్విస్థానిక విధేయ సంకేతం మాత్రమే ఉందనుకుందాం
(వేరే \tetoken{స్థిరసంకేతాలు}{!!{constant}s}, \tetoken{ప్రమేయ సంకేతాలు}{!!{function}s} లేదా \tetoken{విధేయ సంకేతాలు}{!!{predicate}s} ఏవీ లేవు---
అయితే $\eq$ మాత్రం ఉంటుంది). $\Domain{N} = \Nat$, అలాగే
$\Assign{<}{N} = \Setabs{\tuple{n,m}}{n < m}$ అయ్యే నిర్మాణం
$\Struct{N}$ అనుకుందాం. కింది వాటిని నిరూపించండి:
\begin{enumerate}
\item $\{ 0 \}$, $\Struct{N}$లో నిర్వచనీయమైనది;
\item $\{ 1 \}$, $\Struct{N}$లో నిర్వచనీయమైనది;
\item $\{ 2 \}$, $\Struct{N}$లో నిర్వచనీయమైనది;
\item ప్రతి $n \in \Nat$కు సమితి~$\{ n \}$, $\Struct{N}$లో
  నిర్వచనీయమైనది;
\item $\Domain{N}$ యొక్క ప్రతి పరిమిత ఉపసమితి, $\Struct{N}$లో
  నిర్వచనీయమైనది;
\item $\Domain{N}$ యొక్క ప్రతి సహ-పరిమిత ఉపసమితి, $\Struct{N}$లో
  నిర్వచనీయమైనది (ఇక్కడ $X \subseteq \Nat$ అయితే, $\Nat \setminus X$
  పరిమితమైనప్పుడు మాత్రమే ఆ ఉపసమితి సహ-పరిమితమైనది).
\end{enumerate}
\end{prob}


\end{document}
